% ===========================================================================
% CHAPTER 5
% ===========================================================================

\chapter{Certificate Acquisition and Dataset Construction}
\label{chap:certificate-acquisition}

\section{Domain Dataset Preparation}
\label{sec:domain-dataset-preparation}

Certificate acquisition begins with normalized domain inputs. The preparation
process supports ranked domain lists, supplied country-specific lists, and
certificate-derived domain sources. Domain names are converted to lowercase,
validated, deduplicated, and combined into CSV datasets suitable for automated
processing. Alternative certificate-parsing paths support extraction of domain
names from large external certificate collections when such data is available.

Table~\ref{tab:committed-input-counts} summarizes the principal acquisition
inputs included with the project. These values describe input size rather than
the number of unique or successfully retrieved certificates.

\begin{table}[H]
  \centering
  \small
  \setlength{\tabcolsep}{5pt}
  \caption{Certificate acquisition input datasets}
  \label{tab:committed-input-counts}
  \begin{tabularx}{\textwidth}{@{}Y>{\Centering\arraybackslash}p{0.20\textwidth}Y@{}}
    \toprule
    \textbf{Dataset} & \textbf{Records} & \textbf{Purpose} \\
    \midrule
    Global domain dataset & 707,084 & Primary domain input for large-scale certificate acquisition \\
    Pakistan domain dataset & 8,185 & Country-level domain scope used as one supported analytical subset \\
    Pakistan IPv4 prefix set & 720 prefixes & Input for the separate experimental IP-based acquisition workflow \\
    \bottomrule
  \end{tabularx}
\end{table}

The final analytical dataset expands beyond the initial input through
CT-assisted discovery. Dataset size should therefore be interpreted as an
operational input characteristic rather than complete coverage of Internet TLS
deployment.

\section{Domain-Based Certificate Crawling}
\label{sec:domain-crawler-architecture}

The domain-based crawler is a concurrent acquisition service backed by
persistent task states. At startup, it validates its inputs, prepares the
processing state, and records domains as pending tasks. Worker threads
atomically claim tasks so that a normal task is processed by only one worker at
a time. Figure~\ref{fig:domain-crawler-flow} summarizes the workflow.

\begin{figure}[H]
  \centering
  \begin{tikzpicture}[
    font=\scriptsize,
    >=Latex,
    action/.style={draw, rounded corners=2pt, align=center, text width=2.45cm, minimum height=1.0cm, fill=black!5},
    state/.style={draw, rounded corners=2pt, align=center, text width=2.45cm, minimum height=1.0cm, fill=black!9},
    flow/.style={-{Latex[length=2mm]}, line width=0.5pt},
    recovery/.style={-{Latex[length=2mm]}, line width=0.5pt, dashed}
  ]
    \node[action] (preflight) at (-5.8,0) {Validate input and\\supporting services};
    \node[action] (load) at (-2.9,0) {Create pending\\domain tasks};
    \node[state] (claim) at (0,0) {Atomic worker claim\\to processing};
    \node[action] (acquire) at (2.9,0) {TLS retrieval and\\certificate parsing};
    \node[state] (done) at (5.8,0) {Enrich and store;\\mark completed};

    \node[state] (doctor) at (-2.9,-2.0) {Reset stale tasks\\to pending};
    \node[state] (retry) at (0,-3.8) {Eligible retry\\after delay};
    \node[state] (failure) at (2.9,-2.0) {Retrieval or\\parsing failure};
    \node[state] (failed) at (5.8,-3.8) {Record error and\\mark failed};

    \draw[flow] (preflight) -- (load);
    \draw[flow] (load) -- (claim);
    \draw[flow] (claim) -- (acquire);
    \draw[flow] (acquire) -- (done);
    \draw[flow] (acquire) -- (failure);
    \draw[flow] (failure) -- (retry);
    \draw[flow] (failure) -- (failed);
    \draw[recovery] (doctor.north) -- (load.south);
    \draw[recovery] (retry.north) -- (claim.south);
  \end{tikzpicture}
  \caption{Concurrent certificate crawling and recovery workflow}
  \label{fig:domain-crawler-flow}
\end{figure}

For each claimed domain, the crawler opens a TCP connection to port~443,
provides the domain through Server Name Indication, and retrieves the peer
certificate. Hostname and trust validation are disabled during this measurement
step so that certificates that are expired, self-signed, untrusted, or
mismatched can still be observed. The retrieved certificate is converted to
PEM and parsed into structured X.509 and certificate-quality data.

Successful records are enriched with the input domain, acquisition time,
logical scope, and leaf classification. The leaf rule excludes certificates
that are self-subject or explicitly marked as Certificate Authorities. This is
a pragmatic analytical classification rather than complete certification-path
validation.

Persistent states distinguish pending, processing, completed, and failed
tasks. Heartbeat timestamps, configurable retries, stale-task reset, monitoring
statistics, and rescue workers provide stage-level recovery support. These
mechanisms improve operational continuity but do not constitute an exactly-once
processing guarantee.

\section{Experimental IP-Based Acquisition}
\label{sec:experimental-ip-crawler}

The project also includes a separate IP-based certificate crawler for
experimental measurement. Country-specific APNIC allocation data is converted
into IPv4 ranges and CIDR prefixes before host addresses are scheduled through
a worker pool. The crawler supports static and interactive modes, configurable
timeouts, and scans with an IP, a supplied hostname, or no Server Name
Indication.

Outcomes distinguish unavailable port~443 services, TCP and TLS timeouts,
handshake errors, parsing failures, successful certificates, and other errors.
Successful certificates are parsed and stored with the scanned IP and
connection context. This workflow uses a separate dataset and is not combined
with the domain-based analytics dashboard, preventing experimental IP
observations from being interpreted as part of the primary certificate
dataset.

\section{Acquisition Limitations and Ethical Considerations}
\label{sec:acquisition-limitations}

Acquisition results depend on network conditions, DNS, endpoint availability,
port~443 accessibility, TLS behavior, parser execution, and processing
infrastructure. A failed acquisition therefore represents an observed failure
category rather than proof that a domain has no certificate. Similarly,
successful collection does not imply browser trust.

Internet-scale scanning must operate within an authorized and controlled
research environment. Responsible operation requires documented target
authorization, rate limits, opt-out procedures, data-retention rules, and
responsible-disclosure practices. The resulting datasets must also be
interpreted as bounded observations rather than complete or statistically
representative measurements of global TLS deployment.
